% Source: physical PDF pp. 318--323 (printed pp. 295--300). Coverage: complete Section 21.1, from its heading through immediately before Section 21.2; Eqs. (21.1.1)--(21.1.25), citations [2] and [3a], one footnote, no figures or tables. Uncertainty: none; source anomalies are recorded in the wave report.
\subsection{Unitarity Gauge}
\label{sec:21.1}

We saw in Chapter 19 that in a theory with a global symmetry group $G$ that
is spontaneously broken to a subgroup $H$, there is a massless `Goldstone'
boson for every independent broken symmetry, in the sense that the mass
matrix $M^2_{nm}$ of real spinless fields $\phi_n(x)$ has a zero eigenvalue
with eigenvector $\sum_m(t_\alpha)_{nm}v_m$ for each independent broken
symmetry generator $t_\alpha$ of $G$. (We are here considering the case
where these Goldstone bosons are included among the elementary spinless
particles represented by scalar or pseudoscalar fields $\phi_n$ appearing
in the Lagrangian; the more general case will be taken up in Section 21.4.)
We also saw that we could rotate away these Goldstone modes, by subjecting
the fields to a $G$ transformation $\gamma^{-1}(x)$
\begin{equation}
 \widetilde{\phi}_n(x)
 =\sum_m\gamma^{-1}_{nm}(x)\phi_m(x),
 \tag{21.1.1}\label{eq:21.1.1}
\end{equation}
such that the new fields are orthogonal to the Goldstone
directions\footnote{Here we are working with a real reducible or irreducible
representation of the symmetry algebra, for which the matrices $it_\alpha$
are real. The transition to a complex representation is described below.}
\begin{equation}
 0=\sum_{nm}\widetilde{\phi}_n(x)(t_\alpha)_{nm}v_m,
 \tag{21.1.2}\label{eq:21.1.2}
\end{equation}
where $v_n$ is the vacuum expectation value,
$v_n\equiv\langle\phi_n(0)\rangle_{\rm VAC}$. After rotating the fields so
that they satisfy Eq.~\eqref{eq:21.1.2}, the Goldstone boson fields then
reemerged as the spacetime-dependent parameters in $\gamma(x)$. The point
of this procedure was that since the Lagrangian was invariant under the
transformation \eqref{eq:21.1.1} with constant $\gamma(x)$, all dependence
on $\gamma(x)$ dropped out except where $\gamma(x)$ is acted on by
derivatives.

On the other hand, if the Lagrangian is invariant not only under constant
$G$ transformations but also under $G$ transformations that depend on
spacetime position, then the transformation \eqref{eq:21.1.1} is a true
symmetry of the theory, and all dependence on $\gamma(x)$ drops out of the
Lagrangian, so that we can simply replace $\phi_n(x)$ everywhere with
$\widetilde{\phi}_n(x)$. This is a choice of gauge, fixed by imposing the
condition \eqref{eq:21.1.2} on $\widetilde{\phi}(x)$ rather than by imposing
conditions on the gauge fields themselves. (For instance, in the
electrodynamics of a charged scalar field $\phi\equiv\phi_1+i\phi_2$, we
can choose a gauge like Lorentz or Coulomb gauge by imposing conditions
$\partial_\mu A^\mu=0$ or $\boldsymbol{\nabla}\cdot\mathbf A=0$, but we can
also choose a gauge by imposing a condition on $\phi$, as for instance that
$\phi$ be real, or in other words by rotating the two-vector
$\{\operatorname{Re}\phi,\operatorname{Im}\phi\}$ into the 1-direction.)
The gauge defined by Eq.~\eqref{eq:21.1.2} is called \emph{unitarity
gauge},~\hyperref[ch21-ref-3a]{[3a]} because in this gauge it will be
obvious that the theory does not have any degrees of freedom with negative
probability, like timelike gauge bosons. More generally, the unitarity gauge
makes manifest the menu of physical particles of the theory.

Eq.~\eqref{eq:21.1.2} shows that there are no Goldstone boson fields in
unitarity gauge. Since the theory is gauge-invariant this means that there
are no physical Goldstone bosons, whatever gauge we choose. What about the
vector bosons? If the $\phi_n$ are elementary canonically normalized scalar
fields, the Lagrangian will contain a term
\begin{equation}
 \mathcal L_\phi
 =-\frac12\sum_n\left(
 \partial_\mu\widetilde{\phi}_n
 -i\sum_{m,\alpha}t^\alpha_{nm}A_{\alpha\mu}
 \widetilde{\phi}_m\right)^2,
 \tag{21.1.3}\label{eq:21.1.3}
\end{equation}
where $t^\alpha$ runs over all the generators of the gauge group $G$.
(From now on we shall drop tildes, it being understood henceforth in this
section that we are already in unitarity gauge.) We are assuming that the
symmetry $G$ is broken by the vacuum expectation value $v_n$ of $\phi_n$,
so in order to see the nature of the particle spectrum, we define a shifted
field $\phi'$:
\begin{equation}
 \phi_n\equiv v_n+\phi'_n.
 \tag{21.1.4}\label{eq:21.1.4}
\end{equation}
(It is sometimes convenient to take $v_n$ in Eqs.~\eqref{eq:21.1.2} and
\eqref{eq:21.1.4} as the vacuum expectation value in the tree
approximation. In order to generate a useful perturbation theory it is only
necessary that $v_n$ should agree with the true vacuum expectation value in
lowest order.) Expanding Eq.~\eqref{eq:21.1.3} to second order in $\phi'$
and $A$, we have:
\begin{equation}
 \mathcal L_{\phi,\mathrm{QUAD}}
 =-\frac12\sum_n\left(
 \partial_\mu\phi'_n
 -i\sum_{m,\alpha}t^\alpha_{nm}A_{\alpha\mu}v_m\right)^2.
 \tag{21.1.5}\label{eq:21.1.5}
\end{equation}
Using Eq.~\eqref{eq:21.1.2}, we see that the $\phi'$--$A$ cross term
vanishes, yielding
\begin{equation}
 \mathcal L_{\phi,\mathrm{QUAD}}
 =-\frac12\sum_n\partial_\mu\phi'_n\partial^\mu\phi'_n
 -\frac12\sum_{\alpha\beta}\mu^2_{\alpha\beta}
 A_{\alpha\mu}A_\beta^\mu,
 \tag{21.1.6}\label{eq:21.1.6}
\end{equation}
where
\begin{equation}
 \mu^2_{\alpha\beta}
 \equiv-\sum_{nm\ell}t^\alpha_{nm}t^\beta_{n\ell}v_mv_\ell.
 \tag{21.1.7}\label{eq:21.1.7}
\end{equation}
Combining this with the quadratic terms in the Yang--Mills Lagrangian
$-\frac14F_{\alpha\mu\nu}F_\alpha^{\mu\nu}$, we see that the vector
particles have a mass matrix $\mu^2_{\alpha\beta}$. In our notation the
generators $t^\alpha_{nm}$ are proportional to gauge coupling constants, so
Eq.~\eqref{eq:21.1.7} yields vector boson masses that are also proportional
to these coupling constants.

Let's take a look at some of the algebraic properties of
$\mu^2_{\alpha\beta}$. Since $t^\alpha_{nm}$ is imaginary and antisymmetric
(and hence Hermitian) the matrix $\mu^2_{\alpha\beta}$ is real, symmetric,
and positive. Also, if a certain real linear combination of generators
$\sum_\alpha c_\alpha t_\alpha$ is unbroken, then
\begin{equation}
 \sum_{\alpha,m}c_\alpha(t_\alpha)_{nm}v_m=0,
 \tag{21.1.8}\label{eq:21.1.8}
\end{equation}
in which case Eq.~\eqref{eq:21.1.7} shows that, as noted by
Kibble,~\hyperref[ch21-ref-2]{[2]}
\begin{equation}
 \sum_\beta\mu^2_{\alpha\beta}c_\beta=0.
 \tag{21.1.9}\label{eq:21.1.9}
\end{equation}
That is, we still have a massless gauge boson for every unbroken gauge
symmetry. The converse is also true; from Eq.~\eqref{eq:21.1.7} we see that
for arbitrary real constants $c_\alpha$
\begin{equation}
 \sum_{\alpha\beta}\mu^2_{\alpha\beta}c_\alpha c_\beta
 =\sum_n\left(\sum_{\alpha,m}
 c_\alpha\cdot it^\alpha_{nm}v_m\right)^2\geq 0,
 \tag{21.1.10}\label{eq:21.1.10}
\end{equation}
and this can vanish only if $c_\alpha$ satisfies Eq.~\eqref{eq:21.1.8}.

In particular, if there is just one unbroken symmetry
$\sum_\alpha c_\alpha t_\alpha$, then the general gauge field can be written
\begin{equation}
 A_\alpha^\mu=c_\alpha A^\mu+\cdots,
 \tag{21.1.11}\label{eq:21.1.11}
\end{equation}
where `$\cdots$' denotes a linear combination of the gauge fields of
definite non-zero mass, and $c_\alpha$ is the coefficient of $t_\alpha$ in
the single unbroken generator $q$,
\begin{equation}
 q=\sum_\alpha c_\alpha t_\alpha,
 \tag{21.1.12}\label{eq:21.1.12}
\end{equation}
so that the field $A^\mu$ will appear with zero coefficient in the mass
term $-\frac12\sum_{\alpha\beta}\mu^2_{\alpha\beta}
A_\alpha^\mu A_{\mu\beta}$. In order for the coefficient of
$-\frac14(\partial_\mu A_\nu\allowbreak-\partial_\nu A_\mu)^2$ in the kinetic gauge
Lagrangian
$-\frac14\sum_\alpha(\partial_\mu A_{\alpha\nu}
\allowbreak-\partial_\nu A_{\alpha\mu})^2$ to have the canonical value 1, the
$c_\alpha$ must also be normalized
\begin{equation}
 \sum_\alpha c_\alpha^2=1.
 \tag{21.1.13}\label{eq:21.1.13}
\end{equation}
Note that $q$ is the charge to which $A^\mu$ is coupled, in the sense that
\begin{equation}
 \sum_\alpha t_\alpha A_\alpha^\mu=qA^\mu+\cdots,
 \tag{21.1.14}\label{eq:21.1.14}
\end{equation}
where `$\cdots$' again denotes terms involving massive gauge fields. We
shall use these general results in studying the electroweak theory in
Section 21.3.

These results have been derived here for scalar fields that form a real
representation of the gauge group, but they can be straightforwardly
converted to a form appropriate for complex representations. We saw in
Section 19.6 that a complex scalar field $\chi(x)$ that transforms according
to a representation of the gauge group with Hermitian generators $T^\alpha$
may be written as a set of real fields
\begin{equation}
 \phi(x)=
 \begin{pmatrix}
  \operatorname{Re}\chi(x)\\
  \operatorname{Im}\chi(x)
 \end{pmatrix},
 \tag{21.1.15}\label{eq:21.1.15}
\end{equation}
which furnish a real representation of the gauge group with generators
\begin{equation}
 it^\alpha=
 \begin{pmatrix}
  -\operatorname{Im}T^\alpha&-\operatorname{Re}T^\alpha\\
  \operatorname{Re}T^\alpha&-\operatorname{Im}T^\alpha
 \end{pmatrix}.
 \tag{21.1.16}\label{eq:21.1.16}
\end{equation}
Inserting Eqs.~\eqref{eq:21.1.15} and \eqref{eq:21.1.16} in
Eq.~\eqref{eq:21.1.7} gives the vector boson mass matrix
\eqref{eq:21.1.7} as
\begin{equation}
 \mu^2_{\alpha\beta}
 =\operatorname{Re}\left(
 \langle\chi\rangle_{\rm VAC}^\dagger,
 T_\alpha T_\beta\langle\chi\rangle_{\rm VAC}\right)
 =\frac12\left(
 \langle\chi\rangle_{\rm VAC}^\dagger,
 \{T_\alpha,T_\beta\}\langle\chi\rangle_{\rm VAC}\right).
 \tag{21.1.17}\label{eq:21.1.17}
\end{equation}

Now let's take a closer look at the vector field propagator. Including the
quadratic term in the Yang--Mills Lagrangian, the part of the total
Lagrangian that is quadratic in $A$ is
\begin{align}
 &-\frac14\sum_\alpha
 (\partial_\lambda A_{\alpha\nu}-\partial_\nu A_{\alpha\lambda})^2
 -\frac12\sum_{\alpha\beta}\mu^2_{\alpha\beta}
 A_{\alpha\lambda}A_\beta^\lambda
 \notag\\
 &\qquad={}
 \frac12\sum_{\alpha\beta}
 A_\alpha^\nu\mathcal D_{\alpha\nu,\beta\lambda}(\partial)
 A_\beta^\lambda+\text{total derivatives},
 \tag{21.1.18}\label{eq:21.1.18}
\end{align}
where
\begin{equation}
 \mathcal D_{\alpha\nu,\beta\lambda}(\partial)
 =\delta_{\alpha\beta}
 [\eta_{\nu\lambda}\Box-\partial_\nu\partial_\lambda]
 -\mu^2_{\alpha\beta}\eta_{\nu\lambda}.
 \tag{21.1.19}\label{eq:21.1.19}
\end{equation}
Suppose for simplicity that all gauge symmetries are broken, so that
$\mu^2_{\alpha\beta}$ has no zero eigenvalues. According to the general
rules described in Section 9.4, the gauge-field propagator in momentum
space is
\begin{equation}
 \Delta_{\alpha\nu,\beta\lambda}(k)
 =-(\mathcal D^{-1})_{\alpha\nu,\beta\lambda}(ik)
 =\left[(k^2+\mu^2)^{-1}
 \left(\eta_{\nu\lambda}+\mu^{-2}k_\nu k_\lambda\right)
 \right]_{\alpha\beta}.
 \tag{21.1.20}\label{eq:21.1.20}
\end{equation}
Because of the $k_\nu k_\lambda$ term, the propagator has an asymptotic
behavior $\Delta(k)\sim O(k^0)$, which does not allow us to use the usual
power-counting arguments to prove renormalizability. Fortunately, as we
shall see in the next section, there is another gauge in which
renormalizability is obvious, at the price of obscuring the particle content
of the theory.

\begin{center}
$\ast\quad\ast\quad\ast$
\end{center}

It is important to note that although Goldstone bosons suddenly reappear in
the physical spectrum in the limit of zero gauge couplings, physical matrix
elements are perfectly continuous in this limit. This is because in
unitarity gauge the gauge bosons do not entirely decouple for zero gauge
couplings. Consider the matrix element for a scattering process
$A+B\to C+D$, where $A$ and $C$ belong to some representation of the gauge
algebra, and $B$ and $D$ belong to some different representation of the
gauge algebra. The $S$-matrix element for this process receives a
contribution from vector boson exchange
\begin{align}
 S_{CD,BD}={}&i(2\pi)^4
 \delta^4(p_A+p_B-p_C-p_D)
 \bra C J_{N\alpha}^{\nu}\ket A
 \Delta_{\alpha\nu,\beta\lambda}(k)
 \bra D J_{N\beta}^{\lambda}\ket B,
 \tag{21.1.21}\label{eq:21.1.21}
\end{align}
where $k=p_A-p_C=p_D-p_B$, and $J_{N\alpha}^{\nu}$ is the current to which
the gauge bosons are coupled, with the subscript $N$ to remind us that in
this gauge we omit the Goldstone boson pole term in this current. This
current is proportional to gauge coupling constants, so the only terms in
Eq.~\eqref{eq:21.1.21} that survive in the limit of zero gauge couplings are
the ones involving the matrix $\mu^{-2}$, which becomes singular in this
limit. Hence for zero gauge couplings, the $S$-matrix element is
\begin{align}
 S_{CD,BD}\longrightarrow{}&
 i(2\pi)^4\delta^4(p_A+p_B-p_C-p_D)
 k_\nu k_\lambda(\mu^{-2})_{\alpha\beta}
 \bra C J_{N\alpha}^{\nu}\ket A
 \frac1{k^2}
 \bra D J_{N\beta}^{\lambda}\ket B.
 \tag{21.1.22}\label{eq:21.1.22}
\end{align}
The gauge coupling constants in the currents are cancelled by the gauge
coupling constants in the matrix $\mu^2$.

Now let us compare this result with the contribution of Goldstone boson
exchange. For vanishing gauge couplings there is a set of Goldstone bosons
$B_\alpha$ associated with generator $t_\alpha$ (assuming for simplicity
that all gauge symmetries are spontaneously broken) with the kinematic term
in the Lagrangian given by
$-\frac12\partial_\mu\phi_\alpha\partial^\mu\phi_\beta
\left(\sum_nZ_{\alpha n}Z_{\beta n}\right)$, where $Z_{\gamma n}$ is the
component of the Goldstone boson $B_\gamma$ in the spinless field $\phi_n$,
defined by
Eq.~\hyperref[eq:19.2.39]{(19.2.39)}. In order that the Goldstone boson
fields be canonically normalized, we must therefore have
\begin{equation}
 \sum_n Z_{\alpha n}Z_{\beta n}=\delta_{\alpha\beta}.
 \tag{21.1.23}\label{eq:21.1.23}
\end{equation}
According to Eq.~\hyperref[eq:19.2.49]{(19.2.49)}, the Goldstone boson
$B_\alpha$ couples to the currents $J_{N\beta}^{\mu}$ with a coupling
constant $F^{-1}_{\alpha\beta}$, where $F_{\alpha\beta}$ is the coupling of
the Goldstone boson associated with generator $t_\alpha$ to the current
$J_{N\beta}^{\nu}$ defined by Eq.~\hyperref[eq:19.2.38]{(19.2.38)}. Thus
the exchange of Goldstone bosons would give a scattering matrix element
\begin{align}
 S_{CD,BD}={}&i(2\pi)^4\delta^4(p_A+p_B-p_C-p_D)
 k_\nu k_\lambda
 F^{-1}_{\alpha\gamma}F^{-1}_{\beta\gamma}
 \bra C J_{N\alpha}^{\nu}\ket A
 \frac1{k^2}
 \bra D J_{N\beta}^{\lambda}\ket B.
 \tag{21.1.24}\label{eq:21.1.24}
\end{align}
But Eqs.~\hyperref[eq:19.2.40]{(19.2.40)}, \eqref{eq:21.1.7}, and
\eqref{eq:21.1.23} give the vector boson mass matrix as
\begin{equation}
 \mu^2_{\alpha\beta}
 =\sum_n F_{\alpha\gamma}Z_{\gamma n}
 F_{\beta\delta}Z_{\delta n}
 =F_{\alpha\gamma}F_{\beta\gamma},
 \tag{21.1.25}\label{eq:21.1.25}
\end{equation}
so in the limit of vanishing gauge boson coupling, the gauge boson exchange
matrix element \eqref{eq:21.1.22} is the same as the Goldstone boson
exchange matrix element \eqref{eq:21.1.24}. This argument can be reversed;
the requirement of continuity at zero gauge coupling can be used to derive
a formula for the gauge boson masses even in the case where all other
couplings are strong.
